\pageof{One construction, four reasons it works}\label{sec:template}
Pick a predictor \(X\) and a predicted output \(S\). We will engineer an experiment in which reading the \emph{whole announcement} at \(X\) tells us \(S\) perfectly. Which of the old inputs can a classical explanation dispense with? Consider this small example of the construction.
\trio{\(G\): choose \(X,S\)}{\siblingfig{0}}{\(G_1\): restrict \(S\)}{\siblingfig{1}}{\(G'\): forget \(X\)}{\siblingfig{2}}
\captext{In this slice of a larger graph, \(V\) sends an input that \(X\) cannot read. \(G_1\) deletes that arrow. The source shared by \(S,X\) supplies classical data to \(S\). With a childless \(X\) removed, a source serving \(S\) alone becomes private randomness.}
For a sibling predictor, \(S\) keeps its parents seen by \(X\): inputs that also enter \(X\), including shared sources. It loses every other parent. Each retained quantum source can still send a quantum state to its recipients other than \(S\), and sends shared classical data to the whole group. For a visible parent predictor use the simpler rule on the next page.
\quant Start with \emph{any} experiment witnessing a gap in \(G'\). Use the construction for forgetting a node to run it in \(G_1\). For a shared-source predictor, put \(S\)'s random seed in the shared classical data. Then \(X\) can compute \(S\) and announce \((x_{\rm old},s_{\rm prediction})\); its children still use \(x_{\rm old}\). A predicting parent can instead send the seed along its arrow, as on the next page. For several predictors, each reports the classical inputs it sees, so their combined record determines \(S\). This experiment also runs in \(G\), ignoring the deleted inputs.
\begin{rulebox}{What a hypothetical classical explanation must satisfy}
\textbf{P1.} The sources in the classical explanation are independent. Given them, the observed nodes are sampled in causal order, each from its own inputs in \(G\) and private randomness.\\
\textbf{P2.} \(S\) can be read without error from the full \(X\), written \(H(S\mid X)=0\).\\
\textbf{P3.} The observed announcements obey the observable independences of \(G_1\), with the predictor present. The constructed experiment can run there even with the extra field. Observable graphical separations hold for quantum experiments too \cite{HLP}.
\end{rulebox}
\clas We need a classical model of \(G_1\) for these same observations. The next four arguments give sufficient reasons it exists. Delete a childless predictor, or sum over a predictor with children and relay its inputs as on p.~\pageref{sec:forget}. A split predictor loses just its predicting copy. The result explains the original witness for \(G'\) classically, a contradiction. Only this final graph is handed to another piggyback.

\pageof{Fritz I. A parent predicts its child}\label{sec:parent}
\begin{rulebox}{If \(X\) sends its outcome to \(S\), prediction itself suffices}
Keep the arrow \(X\to S\). Delete every other input of \(S\), including its membership in sources shared with \(X\). The same holds for several predicting parents by keeping their arrows.
\end{rulebox}
\duo{A parent can also share a source}{\vtx{x}{0}{1.8}{X}\vtx{s}{3.4}{1.8}{S}\vtx{v}{3.4}{-.3}{V}
 \qsrc{l}{1.4}{.4}{x,s}\de{x}{s}\de[dropedge]{v}{s}}
{The restricted graph, before removal}{\vtx{x}{0}{1.8}{X}\vtx{s}{3.4}{1.8}{S}\vtx{v}{3.4}{-.3}{V}\de{x}{s}
 \node[lab,below] at (1.7,1.65) {send \(x\) and a seed};}
\quant In the restricted experiment \(S\) can sample its output as \(h(x,r)\) using a fresh random seed \(r\). Generate \(r\) at \(X\) and append it to \(x\). The child now outputs exactly the same \(h(x,r)\), and the predictor's full report determines it.
\clas P2 says \(S=g(X)\) on every event of positive probability. In any classical explanation replace \(S\)'s old instruction by this function. It uses a permitted input and preserves the full joint distribution. Here P1 and P2 suffice.

\begin{validbox}{Keep the quantum resource by predicting one field}
Take the chain \(0\to1\to2\to3\), with independent quantum sources on \(01,02,13\) (digits list recipients). Let 0 predict one field \(1'\) of the output at 1. Split 1 first. Before removing 0, the predicted field reads only 0's visible announcement. The original field 1 still receives its share of source \(13\). If we had restricted the entire laboratory 1, that share would have disappeared. The full sequence is drawn on p.~\pageref{sec:chain}.
\end{validbox}
After \(X\) leaves, its children may share classical data from the relay. This explains why shared sources can reappear in the final picture even though \(S\) kept only one arrow in the middle picture.

\pageof{Fritz II. What can independent inputs predict?}\label{sec:separate}
Write \(B\) for the inputs \(S\) keeps and \(A\) for the inputs it loses, including its own private noise. In the original graph, expose each hidden source as a node. Check whether the predictors outside \(B\) are separated from \(A\) once \(B\) is known. The usual path check is called \emph{d-separation}: a path closes at a conditioned ordinary intermediate node, or at a meeting of two arrowheads unless that meeting or a descendant is conditioned on.
\begin{rulebox}{The classical direction in a four-cell picture}
Fix the kept inputs at \(B=b\). Classical instructions give \(S=h(b,A)\), and P2 gives \(S=g(X)\). Graphical separation makes \(A\) and \(X\) independent at this fixed \(b\). For two possible values of each, the common entry would have to satisfy
\[
\begin{array}{c|cc}
 & X=x_0 & X=x_1\\ \hline
 A=a_0 & h(b,a_0)=g(x_0)&h(b,a_0)=g(x_1)\\
 A=a_1 & h(b,a_1)=g(x_0)&h(b,a_1)=g(x_1)
\end{array}
\]
Independence puts every combination of positive-probability row and column on the table. All entries agree. Thus \(S\) can use a function of \(b\) alone. Apply this for each \(b\). The same reasoning applies to larger alphabets, almost surely.
\end{rulebox}
\quant The shared input gives the predictor the data needed to announce \(S\) as on p.~\pageref{sec:template}. The separation check is a fact about the original graph used for the classical argument. That proof needs no observable independence from P3. With several targets, replace one instruction at a time. Each replacement preserves the joint distribution, so the next prediction still holds.
\textbf{Example: the tetrahedron.} There is an independent source for each triple of laboratories and no visible arrows. Laboratory 0 predicts 1. Read each row as one source, each mark as one recipient.
\duo{Original sources}{\cols{0/0,1/1,2/2,3/3}
\members{-.35}{\(012\)}{0,1,2}\members{-1.0}{\(013\)}{0,1,3}
\members{-1.65}{\(023\)}{0,2,3}\members{-2.3}{\(123\)}{1,2,3}}
{After restriction and removing 0}{\cols{1/1,2/2,3/3}
\classmembers{-.35}{classical}{1,2}\classmembers{-1.0}{classical}{1,3}\members{-1.65}{quantum}{2,3}}
\captext{Fix the sources \(012,013\) seen by both 0 and 1. The unseen source \(123\) is independent of the remaining inputs at 0. The new classical links carry the shared data, and 2 and 3 retain a quantum source. This is the triangle seed.}

\pageof{The triangle hides two Bell settings}\label{sec:triangle}
Fritz's original experiment \cite{Fritz} illustrates why predicting two outputs together matters. Let \(A,B\) share an entangled state. A second source supplies a random setting \(x\) to \(A,R\); a third independently supplies \(y\) to \(B,R\). Ask the three laboratories to announce \((x,a)\), \((y,b)\), and \((x,y)\) respectively. Here \(a,b\) are the two measurement outcomes.
\duo{The three laboratories}{\trianglefig{0}}
{Show each setting as its own output}{\trianglefig{1}}
\captext{The right panel shows the restricted experiment. The small laboratories \(x,y\) are the setting fields split from \(A,B\). The record at \(R\) predicts both. The two blue sources are explicitly classical; an all-quantum triangle can run this experiment too.}
\duo{Forget the record at \(R\)}{\trianglefig{2}}
{A Bell test on the surviving fields}{
\node[draw=ccol,rounded corners,fill=ccol!6,align=center] at (2,2.45) {independent choices \(x,y\)};
\node[draw=qcol,rounded corners,fill=qcol!7,align=center] at (2,.95) {measure the shared state};
\node[align=center] at (2,-.15) {read \(a,b\) at settings \(x,y\)};}
\quant Choose any Bell experiment in the final graph. The two settings sources can include the random seeds generating the recorded settings. In the triangle \(R\) receives these seeds and reports both settings. Forgetting \(R\) leaves exactly the Bell distribution.
\clas In a classical triangle, the field \(x\) might initially use the hidden source between \(A\) and \(B\). Given the shared input of \(A,R\), all remaining inputs of \(A\) are independent of the remaining inputs at \(R\). If both always report the same \(x\), the table on p.~\pageref{sec:separate} forces \(x\) to depend only on that shared input. Apply the same argument to \(y\) on the other side. The two settings are generated independently of one another and of the source between \(A,B\). A classical explanation would now give a Bell-local model for the surviving fields, impossible for a Bell-violating experiment.
\begin{validbox}{A gap you can read directly from the triangle records}
Choose independent fair setting bits \(x,y\). Count a win when \(a\oplus b=xy\), where \(\oplus\) is addition modulo two. A classical Bell strategy wins at most three of the four setting pairs on average: for fixed local responses, adding the four left sides modulo two gives 0 (each response appears twice), while adding the four right sides gives 1. Shared random choices only average these bounds. On the state \((|00\rangle+|11\rangle)/\sqrt2\), measuring the Pauli observables \(Z\) or \(X\) at \(A\) and \((Z+X)/\sqrt2\) or \((Z-X)/\sqrt2\) at \(B\) wins with probability \((2+\sqrt2)/4>3/4\). The eigenvalues \(+1,-1\) are recorded as \(0,1\). Laboratory \(R\)'s two predictions certify the setting fields used to score the game.
\end{validbox}

\pageof{Fritz III. What the extra independence buys}\label{sec:kpc}
Khanna, Pusey and Colbeck (KPC) \cite{KPC} give a graph where the last path test fails. Use their names: \(C,D,E,F\) are observed. Source \(A\) feeds \(E,F\) and source \(B\) feeds \(D,F\). Ask \(F\) to predict \(E\). The restricted \(E\) keeps only the source shared with \(F\). Retain the ordinary output \(F\) by forgetting just its predicting field.
\trio{Original graph \(G\)}{\kpcfig{0}}{Restricted graph \(G_1\)}{\kpcfig{1}}{Classical explanation}{\kpcfig{2}}
\captext{Left: fixing \(A\) leaves \(F\leftarrow B\to D\) open. Middle: removing the two dashed arrows gives \(E\ind(C,D)\), while \(D,F\) may still measure a shared quantum state. Right: if the original observations had a classical explanation, the proof below would give this model with independent classical sources \(A',B\).}
\quant Begin with that Bell experiment: \(C\) supplies a setting to \(D\), the shared classical source supplies the setting \(E\) to \(F\), and \(D,F\) share a quantum state. Give \(F\) a copy of the data determining \(E\), including its random seed. In the original graph the extra arrows into \(E\) can be ignored. This realizes P2 and the observable independence in P3. In the deductions below \(F\) denotes the full report including its extra field.
\clas Write \(U=(C,D)\). In a proposed classical explanation, \(J\ind K\mid L\) means that once \(L\) is known, learning \(J\) gives no further information about \(K\).
\begin{enumerate}
\item \(E\ind B\mid U\) follows in \(G\): given \(C,D\), \(E\) uses only the independent source \(A\) and its own noise. P3 gives \(E\ind U\). Combining these yields \(E\ind(B,U)\).
\item \(F\ind U\mid B\) follows in \(G\): fixing \(B\) separates the inputs generating \(F\) from \(C,D\). P2 gives \(E\ind U\mid B,F\). Combine the statements to get \((F,E)\ind U\mid B\), then condition on \(E\) to get \(F\ind U\mid B,E\).
\item \(B\ind C\) in \(G\). The two preceding lines now give the factorization
\[
p(b,c,d,e,f)=p(b)\,p(c)\,p(d\mid b,c)\,p(e)\,p(f\mid b,e).
\]
\end{enumerate}
Sample \(A'\) independently with distribution \(p(e)\), let \(E=A'\), and give \(F\) the inputs \(A',B\). The factorization proves these instructions preserve all observed statistics. Forget \(F\)'s extra field. This is the desired classical explanation of the Bell experiment and completes the contradiction.

\pageof{Why the old source can be the wrong one}\label{sec:fresh}
Let \(A,B,C\) be independent fair bits. The symbol \(\oplus\) means addition modulo two. Keep the KPC arrows and give the laboratories the following instructions:
\[
 D=C\oplus B,\qquad E=A\oplus C\oplus D=A\oplus B,\qquad F=(A,B).
\]
Laboratory \(E\) can use the first expression because it receives \(A,C,D\). Laboratory \(F\) can predict it by the second. For each fixed \(C,D\), the fresh fair bit \(A\) makes \(E\) fair, so \(E\ind(C,D)\). Could \(E\) simply toss a private coin after reading the old \(A\)? Given \(A\), \(E\) and \(B\) are perfectly correlated or anticorrelated: \(E=A\oplus B\). A rule using only \(A\) and independent private noise would give \(E\ind B\mid A\). Thus it cannot preserve this classical joint distribution. The table constructs a different joint distribution with the same observations.
\duo{Sample \(A,B,C\) independently}{
\node[draw=qcol,fill=qcol!6,rounded corners,align=left] at (1,1) {\(\begin{array}{cc|cc} A&B&E&F\\\hline 0&0&0&00\\0&1&1&01\\1&0&1&10\\1&1&0&11\end{array}\)};}
{Sample \(A',B,C\) independently}{
\node[draw=ccol,fill=ccol!6,rounded corners,align=left] at (1,1) {\(\begin{array}{cc|cc} A'&B&E&F\\\hline 0&0&0&00\\1&1&1&01\\1&0&1&10\\0&1&0&11\end{array}\)};}
\captext{Each row has probability \(1/4\) and is repeated for \(C=0,1\) with equal weights. Corresponding rows have identical \(B,E,F\), hence identical \(C,D,E,F\). On the right use \(E=A'\), \(D=C\oplus B\), and \(F=(A'\oplus B,B)\).}
\begin{rulebox}{Two sufficient ways to prove a restricted graph has a classical explanation}
\textbf{Keep the old sources.} Use P1, P2 and P3 to deduce the local independence statements of the restricted graph. They say each node can be sampled given its permitted parents. This test can certify restrictions directly.

\textbf{Replace one source.} When one predicted node \(S\) keeps exactly one parent, a hidden source \(L\), remove \(L\) on paper. Make \(S\) a root and send its value to every other former recipient of \(L\). Prove the local independence statements of this graph over the remaining variables. Then a fresh independent source carrying \(S\) realizes the restricted graph, exactly as on p.~\pageref{sec:kpc}.
\end{rulebox}
\textbf{How the independence test works.} Start with graphical independences of the original classical model, the observable ones from the restriction, and the fact that once the predictor is given, learning other variables adds no information about \(S\). This last fact continues to hold with extra variables conditioned on. Close the list under symmetry, dropping variables from one side, moving part of that side into the conditioning set, and the reverse combining step (contraction). These are the \emph{semigraphoid rules}. The proof on the previous page is a short instance. This calculation uses a consequence of perfect prediction; independence from other variables given \(X\) alone also holds for an unrelated fresh random bit.

\pageof{Fritz IV. Read the same proof in bits}\label{sec:entropy}
Read \(H(E\mid F)=0\) as zero uncertainty about \(E\) once \(F\) is known. Read \(I(J;K\mid L)=0\) as independence of \(J,K\) given \(L\). Both quantities are nonnegative for the joint distribution of a classical explanation.
\begin{rulebox}{An entropy calculation for the KPC graph}
Again write \(U=(C,D)\). P1 supplies \(I(E;B\mid U)=0\) and \(I(F;U\mid B)=0\). P3 supplies \(I(E;U)=0\). The chain rule gives
\[
 I(E;B,U)= I(E;U)+I(E;B\mid U)=0.
\]
P2 supplies \(H(E\mid F)=0\). Consequently
\[
 0\le I(E;U\mid B,F)\le H(E\mid B,F)\le H(E\mid F)=0.
\]
Use the chain rule in its two orders:
\[
 \underbrace{I(F;U\mid B)+I(E;U\mid B,F)}_{0}
     = I(E;U\mid B)+I(F;U\mid B,E).
\]
Both terms on the right must vanish. These are the same ingredients that allow the fresh-source factorization on p.~\pageref{sec:kpc}.
\end{rulebox}
\quant The experiment on p.~\pageref{sec:kpc} again lets \(F\) announce \(E\) exactly.
\clas For the general test, take one entropy coordinate for every nonempty subset of the classical variables. Impose the Shannon inequalities (nonnegative conditional entropies and conditional mutual informations), the original graph's local independence equalities, full \(H(S\mid X)=0\), and the visible independence equalities of \(G_1\). A local equality for node \(v\) says \(I(v;N\mid\Pa(v))=0\), where \(N\) contains its non-descendants other than itself and its parents.
\textbf{Why one linear program suffices for a proposed conclusion.} Sum the local quantities for the chosen graph, calling the sum \(T\). Each summand is nonnegative. To test if all vanish, ask whether the constraints permit \(T\ge1\). All constraints before the last are homogeneous, so any positive \(T\) could be scaled to 1. Infeasibility forces \(T=0\) for every compatible classical model of \(G\). Each required independence then holds, so the construction on p.~\pageref{sec:fresh} gives a classical explanation in \(G_1\). The solver can certify \(T\le0\) by combining the premise equalities and nonnegative multiples of the inequalities. One solve checks each proposed set of conclusions; testing both choices on p.~\pageref{sec:fresh} can require two. Feasibility leaves the question undecided because a Shannon vector need not come from a distribution.

\begin{ideabox}{Why agreement on the census does not make the two tests equivalent}
For four arbitrary variables, expand each information quantity into entropies. The following two sums agree term for term after cancellation:
\begin{align*}
 &I(A;B\mid C,D)+I(C;D\mid A)+I(C;D\mid B)+I(A;B)\\
={}&I(C;D\mid A,B)+I(A;B\mid C)+I(A;B\mid D)+I(C;D).
\end{align*}
If the first four terms vanish, the last four vanish too. The semigraphoid rules applied to those premises derive none of the four conclusions \cite{Repo}.
\end{ideabox}
